%% 第二章

\chapter{图表示例}
\chaptermark{图表示例}

\section{表格}
如\autoref{tab:undergraduate-example} 所示，该表展示了两个变量的加法与乘法结果，可作为表格排版与引用的示例。
\begin{table}[!h]
    \caption{表格标题}
    \label{tab:undergraduate-example}
    \begin{tabular}{|c|c|c|c|}
        \hline
        $A$ & $B$ & $A+B$ & $A\times B$ \\ \hline
        1   & 6   & 7     & 6           \\ \hline
        2   & 7   & 9     & 14          \\ \hline
        3   & 8   & 11    & 24          \\ \hline
    \end{tabular}
\end{table}

\section{插图}
测试下文章内的图片引用。如\autoref{fig:undergraduate-example} 和\autoref{fig:undergraduate-example2} 所示，
这是两幅插图。在这其中\autoref{subfig:undergraduate-example2-subfig1} 是第一幅子图，
\autoref{subfig:undergraduate-example2-subfig2} 是第二幅子图。

\begin{figure}[htb]
    \includegraphics[scale=0.5]{example.png}
    \caption{这里是个普通的标题}
    \label{fig:undergraduate-example}
\end{figure}

\begin{figure}[htb]
    \begin{subfigure}[t]{0.49\textwidth}\centering
        \includegraphics[scale=0.5]{example.png}
        \caption{这是第一个子图标题}
        \label{subfig:undergraduate-example2-subfig1}
    \end{subfigure}
    \hfill
    \begin{subfigure}[t]{0.49\textwidth}\centering
        \includegraphics[scale=0.5]{example.png}
        \caption{这是第二个子图标题}
        \label{subfig:undergraduate-example2-subfig2}
    \end{subfigure}
    \caption{这里是另一个普通的标题}
    \label{fig:undergraduate-example2}
\end{figure}

\section{公式示例}
这里给出正文中的公式引用示例。行内公式可写为 $E=mc^2$，独立编号公式如\autoref{eq:undergraduate-quadratic-formula} 所示，多行推导公式如\autoref{eq:undergraduate-gaussian-sum} 所示。

\begin{equation}
    x = \frac{-b \pm \sqrt{b^2 - 4ac}}{2a}
    \label{eq:undergraduate-quadratic-formula}
\end{equation}

\begin{equation}
    \begin{aligned}
        \sum_{i=1}^{n} i &= 1 + 2 + \cdots + n \\
                         &= \frac{n(n+1)}{2}
    \end{aligned}
    \label{eq:undergraduate-gaussian-sum}
\end{equation}
